\documentclass[10pt,a4paper]{article}

\usepackage[a4paper,margin=2.3cm,top=2.1cm,bottom=2.1cm]{geometry}
\usepackage{amsmath}
\usepackage{fontspec}
\usepackage{unicode-math}
\setmainfont{texgyrepagella}[Extension=.otf,UprightFont=*-regular,BoldFont=*-bold,ItalicFont=*-italic,BoldItalicFont=*-bolditalic]
\setsansfont{texgyreheros}[Extension=.otf,UprightFont=*-regular,BoldFont=*-bold,ItalicFont=*-italic,BoldItalicFont=*-bolditalic,Scale=0.92]
\setmathfont{texgyrepagella-math.otf}
\usepackage[croatian]{babel}
\usepackage{microtype}
\usepackage{xcolor}
\usepackage{booktabs,tabularx,array}
\usepackage{enumitem}
\usepackage{titlesec}
\usepackage[hidelinks]{hyperref}

\definecolor{accent}{HTML}{0F5E63}
\definecolor{muted}{HTML}{55606A}

\titleformat{\section}{\sffamily\large\bfseries\color{accent}}{\thesection.}{0.5em}{}
\titleformat{\subsection}{\sffamily\normalsize\bfseries}{\thesubsection}{0.5em}{}
\titlespacing*{\section}{0pt}{1.3em}{0.45em}
\titlespacing*{\subsection}{0pt}{0.9em}{0.25em}

\setlength{\parindent}{0pt}
\setlength{\parskip}{0.45em}
\setlist{nosep,leftmargin=1.3em,topsep=0.2em}
\newcolumntype{Y}{>{\raggedright\arraybackslash}X}
\renewcommand{\arraystretch}{1.25}
\emergencystretch=1em
\widowpenalty=10000
\clubpenalty=10000

\newcommand{\R}{\mathbb{R}}

\hypersetup{pdftitle={NMF seminar: prijedlog tema i način rada},pdfauthor={Lovro}}

\begin{document}

{\sffamily\LARGE\bfseries\color{accent} NMF seminar: prijedlog tema i način rada\par}
\vspace{0.3em}
{\small\color{muted} Matrične i tenzorske metode u analizi podataka \quad·\quad Lovro, za Saru i Mislava \quad·\quad 1.~studenoga 2026.\par}
\vspace{0.6em}

Moramo detaljno obraditi nenegativnu matričnu faktorizaciju, implementirati je u Pythonu i
primijeniti na nov konkretan problem. Ovdje su četiri moguće izvedbe. Svaka ima drugi problem i
drugi algoritam iz literature, pa izbor teme određuje i matematiku koju ćemo braniti.
\textbf{Predlažem temu~1, a temu~2 kao rezervu.} Ništa još nije odlučeno.

Podaci i licence provjereni su 1.~studenoga 2026., ali nijedna metoda još nije pokrenuta, pa su
očekivani rezultati zasad samo pretpostavke. Svuda koristimo model
$A\approx WH$, gdje je $A\in\R_{\geq0}^{m\times n}$, $W\in\R_{\geq0}^{m\times r}$ i
$H\in\R_{\geq0}^{r\times n}$.

\section{Četiri teme}

\subsection*{Tema 1: Kartiranje močvare Kopački rit iz satelitskih snimaka}
Pitanje je može li se bez ručno označenih podataka iz snimaka Sentinel-2 odvojiti vodu, močvarnu
vegetaciju, šumu i polja te izmjeriti kako se površina pod vodom mijenja kroz godinu.
Stupci matrice $A$ su pikseli, a retci spektralni kanali na nekoliko datuma.
Spektar piksela smatramo kombinacijom spektara materijala, što je linearni model miješanja iz
Gillisa (odj.~2.3): stupci $W$ su spektri materijala, a stupci $H$ njihovi udjeli.
Najčišće piksele tražimo algoritmom SPA (Gillis, odj.~3.2.2) i njima pokrećemo faktorizaciju ranga~2.
Svaki čvor dijelimo pomoću ANLS-a sa zatvorenom formulom (L6, odj.~3), pa nastaje stablo materijala.
Završne karte udjela dobivamo nenegativnim najmanjim kvadratima uz fiksirani $W$.
Rezultat uspoređujemo s k-means i pragom indeksa NDWI, a slaganje mjerimo s kartom ESA WorldCover 2021,
koja je i sama nesavršena referenca (ukupna točnost 76,7~\%).
Jedna pločica (34TCR) pokriva cijelo područje i za 2025.\ ima 39 snimaka s manje od 30~\% oblaka na razini pločice.
Rizici su oblaci, malo proljetnih snimaka i pravilo dijeljenja mNDCG iz L6, koje je smišljeno za
riječi pa će ga možda trebati zamijeniti.

\subsection*{Tema 2: Automatsko prepisivanje gitare u note}
Pitanje je može li NMF iz snimke akustične gitare prepoznati koje tonove glazbenik svira i kada.
Stupci $A$ su vremenski okviri, a retci frekvencijski binovi amplitudnog spektrograma.
Spektar u nekom trenutku približno je zbroj spektara tonova koji tada zvone, pa su stupci $W$
predlošci tonova, a retci $H$ njihove aktivacije kroz vrijeme.
Koristimo generaliziranu KL divergenciju i multiplikativna ažuriranja iz L3, gdje je cijeli dokaz
nerasta funkcije cilja pomoću pomoćne funkcije (teorem~2, odj.~6).
L3 tvrdi i konvergenciju prema lokalnom minimumu, ali L4 (odj.~3.1) pokazuje da taj dokaz daje
samo nerast, pa u seminaru razdvajamo što je dokazano, a što nije.
Predloške učimo na pet gitarista, a transkripciju provjeravamo na šestom, i to ponavljamo za svakoga.
Uz fiksirani $W$ problem po $H$ je konveksan, pa je obrada nove snimke brza i lako objašnjiva.
Skup GuitarSet (CC~BY~4.0) ima 360 snimaka šest gitarista i 62.476 označenih tonova, a dovoljna je
arhiva od 657~MB.
Rezultat mjerimo standardnim F1 na razini nota i uspoređujemo ga s Frobeniusovom funkcijom cilja i
jednostavnim odabirom spektralnih vrhova.
Rizik je što isti ton zvuči različito na različitim žicama, a akordi se preklapaju, pa rezultat
može biti skroman.
Dodatno moramo objasniti kratkotrajnu Fourierovu transformaciju.

\subsection*{Tema 3: Predviđanje šahovskog repertoara}
Pitanje je može li se iz malog broja partija predvidjeti koja će otvaranja šahist igrati, što
pomaže u pripremi protiv protivnika.
Stupci $A$ su igrači, a retci otvaranja odvojeno za bijele i crne figure, što je izravna analogija
matrice riječi i dokumenata iz L5.
Zato koristimo težine TF-IDF i HierNMF2 iz L6, pa stablo može otkriti repertoarske stilove igrača.
Težine novog igrača računamo nenegativnim najmanjim kvadratima uz fiksirani $W$, što bi trebalo
raditi i uz malo partija.
Partije dijelimo vremenski: prvi dio mjeseca služi za učenje, a drugi za provjeru predviđanja.
Uspoređujemo se s popularnošću otvaranja i igračevim vlastitim histogramom, a stablo s ECO
klasifikacijom otvaranja.
Baza Lichess je pod licencom CC0, a za srpanj 2014.\ ima 1.048.440 partija, 2.619 naziva otvaranja
i 5.481 igrača s barem 100 partija.
Ta datoteka ima oko 200~MB, pa je obrada izvediva na običnom računalu.
Rizik je što otvaranje ovisi i o protivniku, pa značajke treba pažljivo definirati.
NMF možda neće pobijediti vlastiti histogram, što bi bio legitiman negativan rezultat.

\subsection*{Tema 4: Razdvajanje bojila na histološkim slikama}
Preparati tkiva bojani hematoksilinom i eozinom razlikuju se u boji među laboratorijima i
skenerima, što smeta automatskoj analizi.
Pitanje je može li NMF ranga~2 razdvojiti ta dva bojila i time ujednačiti boje slika.
Nakon Beer-Lambertove transformacije u optičku gustoću piksel je približno nenegativna kombinacija
boja dvaju bojila, pa rang određuje fizika, a ne naš izbor.
Stupci $A$ su pikseli jedne slike, retci tri RGB kanala, stupci $W$ boje bojila, a retci $H$ njihove
gustoće.
Budući da $A$ ima samo tri retka, nejedinstvenost NMF-a možemo nacrtati kao konus u trodimenzionalnom
prostoru.
Prema Vahadaneu i suradnicima dodajemo kaznu $\ell_1$ na $H$, a problem rješavamo HALS-om s mekim
pragom (Gillis, odj.~3.1.5).
Metodu uspoređujemo s Macenkovom, a korist mjerimo jednostavnim klasifikatorom tkiva provjerenim na
drugim pacijentima.
Skup NCT-CRC-HE-100K (CC~BY~4.0) ima 100.000 nenormaliziranih isječaka devet vrsta tkiva i 7.180
isječaka od 50 drugih pacijenata.
Puna arhiva ima 11,7~GB, ali dovoljan je uzorak pojedinačnih slika.
Rizici su potrebno medicinsko predznanje, moguća prethodna normalizacija validacijskih slika i to
što nazivi datoteka ne otkrivaju s kojeg preparata isječak potječe.

\section{Usporedba}

{\small
\begin{tabularx}{\linewidth}{@{}>{\raggedright\arraybackslash}p{2.05cm}YYY>{\raggedright\arraybackslash}p{2.75cm}@{}}
\toprule
\textbf{Tema} & \textbf{Funkcija cilja i algoritam} & \textbf{Veza s literaturom} & \textbf{Vrednovanje} & \textbf{Licenca podataka}\\
\midrule
1. Kopački rit & Frobeniusova; SPA, rang-2 ANLS, stablo & L1 (2.3, 3.1.4, 3.2.2), L6 & Slaganje s WorldCoverom, k-means, NDWI & Sentinel otvoren; WorldCover CC BY 4.0\\
2. Gitara & KL; multiplikativna ažuriranja & L3 (teorem 2), L1 (3.1.2), L4 (3.1) & F1 na notama, provjera na novom gitaristu & CC BY 4.0\\
3. Šah & Frobeniusova; HierNMF2 na TF-IDF & L5, L6 & Predviđanje na kasnijim partijama & CC0\\
4. Histologija & Frobeniusova s $\ell_1$; HALS & L1 (3.1.5), Vahadane i sur. & Usporedba s Macenkom, klasifikator tkiva & CC BY 4.0\\
\bottomrule
\end{tabularx}
}

\textbf{Preporuka.} Tema~1 najviše pokriva glavni obvezni izvor: model miješanja, ANLS i SPA s
dokazom, a stablo dolazi iz L6. Daje i karte koje se dobro prikazuju, ali traži najviše rada s
pripremom podataka. Tema~2 ima najčišći dokaz i najjasniju mjeru uspjeha, tema~3 najbliža je radovima
o dokumentima, a tema~4 ima najzorniju geometriju, ali traži medicinski kontekst.

\section{Predložena metodologija rada: niti}

Posao ne dijelimo na teoriju, kod i primjenu. Dijelimo ga na tri isprepletene niti, a svaka ide od
matematike preko koda i podataka do rezultata te svog dijela seminara i prezentacije. Metoda je ista
za svaku temu.

{\small
\begin{tabularx}{\linewidth}{@{}>{\raggedright\arraybackslash}p{3.2cm}Y>{\raggedright\arraybackslash}p{4.1cm}@{}}
\toprule
\textbf{Nit} & \textbf{Što obuhvaća} & \textbf{Završava}\\
\midrule
1. Od podataka do prve faktorizacije & Što znači $A$ i zašto je nenegativna, funkcija cilja, uvjeti optimalnosti, nejedinstvenost, izvod ažuriranja, prvi kod i mali primjer koji se može provjeriti ručno & Prvom faktorizacijom na stvarnim podacima i tumačenjem faktora\\
2. Od faktorizacije do strukture & Izbor ranga ili stabla, inicijalizacija, zaustavljanje, što je o konvergenciji stvarno dokazano, stabilnost kroz sjemena, usporedba s referentnom metodom & Stabilnom strukturom u podacima i usporedbom\\
3. Od strukture do odgovora & Pretvaranje faktora u odgovor, vrednovanje, nesigurnost, neuspjesi i ograničenja, ponovljivost pokretanja & Odgovorom na pitanje i završnim slikama\\
\bottomrule
\end{tabularx}
}

Primjer za temu~1: nit~1 završava prvom podjelom scene na vodu i kopno, nit~2 stablom materijala, a
nit~3 kartama udjela i krivuljom vode kroz godinu.

\textbf{Kako radimo:}
\begin{itemize}
\item Svaka nit ima nositelja, a ne vlasnika. Nositelj prati nit, a niti možemo mijenjati usput.
\item AI alat piše kod i prve nacrte, što je na seminaru dopušteno. Mi sve provjeravamo, pokrećemo i odgovaramo za svaki redak i tvrdnju.
\item Seminar i prezentacija prate niti, pa svaki dio zajedno ima matematiku, kod i sliku.
\item Imamo tri zajednička prolaza: nakon prve faktorizacije, nakon strukture i prije obrane. Na svakom nit objašnjava netko tko joj nije nositelj.
\item Svatko mora moći objasniti cijeli projekt, pa na probnoj obrani dobiva pitanja iz tuđih niti.
\item Radimo u malim zadacima, svaki na svojoj grani, a PR odobrava drugi član tima. Pitanja i hipoteze zapisujemo prije završnih pokretanja.
\end{itemize}

\textbf{Redoslijed:} izbor teme, nit~1, nit~2, nit~3, završni seminar i prezentacija, zatim ponovno
pokretanje svih rezultata i probna obrana. Rok još ne znamo, pa ovo nije vremenski plan.

\textbf{Trebamo odlučiti:} koju temu biramo, tko je u početku nositelj koje niti te koji su rok i format predaje.

\section*{Izvori}
{\footnotesize\raggedright
\textbf{Literatura seminara.}
L1: N.~Gillis, \emph{The Why and How of Nonnegative Matrix Factorization}, arXiv:1401.5226v2.
L3: D.~D.~Lee, H.~S.~Seung, \emph{Algorithms for Non-negative Matrix Factorization}.
L4: M.~W.~Berry, M.~Browne, A.~N.~Langville, V.~P.~Pauca, R.~J.~Plemmons, \emph{Algorithms and Applications for Approximate Nonnegative Matrix Factorization}.
L5: W.~Xu, X.~Liu, Y.~Gong, \emph{Document Clustering Based On Non-negative Matrix Factorization}, SIGIR 2003.
L6: D.~Kuang, B.~Drake, H.~Park, \emph{Fast Clustering and Topic Modeling Based on Rank-2 Nonnegative Matrix Factorization}, 2015.

\textbf{Dodatno.}
A.~Vahadane i sur., \emph{Structure-Preserving Color Normalization and Sparse Stain Separation for Histological Images}, IEEE TMI, 2016.
M.~Macenko i sur., ISBI 2009, doi:10.1109/ISBI.2009.5193250.
Q.~Xi i sur., \emph{GuitarSet: A Dataset for Guitar Transcription}, ISMIR 2018.

\textbf{Podaci.}
Sentinel-2 L2A: \url{https://earth-search.aws.element84.com/v1};
ESA WorldCover: \url{https://zenodo.org/records/7254221};
GuitarSet: \url{https://zenodo.org/records/3371780};
Lichess: \url{https://database.lichess.org/};
NCT-CRC-HE-100K: \url{https://zenodo.org/records/1214456}.
\par}

\end{document}
